\documentclass[fontsize=14pt]{mybib}

\setheaderlogo{mnr-green}

% ensure scrlayer-scrpage has sufficient footheight
\renewcommand{\DFA}{12mm}
\renewcommand{\DZA}{1.2}
\setlength{\footheight}{30.4pt}
\author{OTS}

\begin{document}
    \section*{Bericht zur Evangelisation Bahnhof Brig}
    \subsection*{Samstag}
    Im Frühjahr 2026, hat mich Fredy darauf aufmerksam gemacht, dass Anton Saljuk im MNR Evangelisation und verteilen von Bibeln durchführt. Ich solle ihn doch mal ansprechen und Frage, ob er so etwas auch in Naters durchführen würde. Anton und ich waren im gleichen EBTC-Kurs \enquote{Bibellesen mit Gewinn} und ich habe ihn dann am letzten Tag angesprochen und gefragt, ob er bei uns in Naters auch mal so ein Bibelverteilen durchführen wollte. Anton hat dann auch direkt zugesagt.\\
    Im Sommer haben wir dann das Datum 10. 10. 2026 ein Samstag reserviert. Anton konnte noch vier Personen zusätzlich organiseren. Elisabeth, Francesco, Gjergi, Jonathan kamen zusammen mit zwei Autos aus Dübendorf und spontan hat sich noch Daniel spontan für diese Aktion gemeldet. Daniel haben wir dann später am Bahnhof Brig getroffen.\\
    Erwartet wurde die Gruppe um 12:00Uhr, gekommen sind sie dann gegen 13:00Uhr. Nicht weil soviel Verkehr war, sondern weil das Wallis so schön ist, dass sie direkt den Furkapass noch mitgenommen haben.
    
    Nachdem wir uns begrüsst, vorgestellt und gemeinsam Gott gedankt haben für die unfallfreie Fahrt sind wir zusammen auf den Bahnhof und dort haben wir gemeinsam die Kisten ausgepackt. Anton hat rund 400 Bibeln, Neues Testament mit Psalmen, mitgenommen. Zusätzlich noch Kinderbibeln, und Bibeln in Italienisch und Französisch.

    Bahnhof Brig ist nicht mehr der grosse Umschlagplatz, sondern ein Endbahnhof und ein Startbahnhof nach Italien und in die \enquote{Üssuschwiiz}. Die Menschen kamen Stossweise und viele waren Feriengäste. Francesco konnte als Italiener direkt die ansprechen. Es gab auch ein paar Asiaten und englisch Sprechende aber leider hatten wir keine englischen Bibeln dabei. Ich war doch überrascht wie viele Jungen Leute eine Bibel genommen haben. Es wurde auch oft gefragt, um welche Übersetzung es sich dabei handelt und ob wir eine Sekte sind.

    Wir haben uns in zwei Gruppen aufgeteilt und so von 13:15Uhr bis 17:20Uhr am Bahnhof Brig Bibeln verteilt. Es wurden rund 200 Bibeln verteilt. Was daraus wird, liegt nun nicht mehr in unseren Verantwortung. Ist das nicht fantastisch? Jesus hat uns den Autrag geben, nur die gute Botschaft zu verbreiten. Was damit dann passiert, darum brauchen wir uns nicht drum zu bekümmern. Wie schlimm wäre das, wenn wir auch für die Bekehrung verantwortlich sind?
    Nach der Rede von Apostel Petrus in Jerusalem, haben sich 3000 Leute bekehrt. Wegen seiner extrem guten Rede? Nein, weil der Heilige Geist die Bekehrung der Menschen übernommen haben. Wir müssen nicht wenn wir eines Tages vor Jesus stehen, ihm eine Liste von Namen vorweisen die wir bekehrt haben. Aber es kann sein, dass Jesus diesen jungen Menschen die sich so für die Verbreitung Evangelium einsetzen, ein Liste von Jesus unter die Nase bekommen, mit den Worten: \enquote{Du guter treuer Knecht!} Lasst uns einfach das tun was Knechte tun. Dem Chef gehorchen. Aber ich schweife ab.

    Das Verteilen der Bibeln war aber nicht die einzige Liebestat, die an dem Wochenende gelaufen ist. Im Hintergrund haben die lieben Schwestern Marie-Therese und Josette für das leibliche Wohl der Gäste gesorgt. Nach dem Verteilen waren wir alle, in unserem Lokal zu einem gemeinsamen Abendessen eingeladen. Es war fantastisch. Es gab eine typische Walliser Käse- und Fleischplatte mit Kartoffel- und Grünem Salat.\\

    Natürlich wurden am Tisch die Erlebnisse vom Nachmittag besprochen und Francesco hat uns aufgeklärt wie man als Christ mit Muslimen umgeht und vor evangelisierenden Muslimen gewarnt. Dann wurden die Gruppen wieder aufgeteilt 3 zum Schlafen nach Birgisch und 2 kamen in Naters bei der Familie Perren unter.

    \subsection*{Sonntag}
    Es zeigte sich, dass der Sonntag ein wunderschöner Herbsttag wird. Die ganze Gruppe war zum Gemeinsamen Mittagessen eingeladen, die wieder von Marie-Therese, Josette und Armin vorbereitet und gekocht wurde. Es gab leckeren Reis mit Geschnetzeltem. Am Sonntag kamen dann auch noch Ueli und Ursula dazu, die leider am Samstag keine Zeit hatten.

    Für die Predigt am Sonntag war Elia Morries eingeteilt. Er kam schon früh 13:13Uhr in Brig an. Das war sehr schön. So gab es noch gute Gespräche und die Gruppe ist dann unter der Leitung von Ursula, vor dem Gottesdienst noch durch Naters spaziert. Janina, Ueli und ich haben den Gottesdienst vorbereitet.

    Da sich die Gruppe nicht verlaufen hat, konnte pünktlich 14:45 Uhr mit dem Gottesdienst gestartet werden. Die von Elias war über Zacharias 3 1-10. Das Ablegen der Sündigen Kleider und sich mit neuen weissen einkleiden lassen. Wichtig sich einkleiden lassen. Jesus ist für meine Sünden gestorben und wenn ich dieses Opfer von ihm von ganzem Herzen annehme, dann wäscht er mich mit seinem Blut rein. Ist das nicht herrlich?

    Nach dem Gottesdienst gab es Kaffee und Kuchen und dann musste Abschied genommen werden. Nicht aber ohne zu versprechen das ganze zu wiederholen.

    \subsection*{Mein Fazit}
    Dieses Wochenende hat mir gezeigt, dass wir auf dem richtigen Weg. Diese gelebte Gemeinschaft und gegenseitige Hilfe und Unterstützung war überwältigend.\\
    Das Mitwirken Gottes war in jedem Dienst spürbar.

    Ich möchte der Gruppe von ganzem Herzen Danken für ihren Dienst. Ich möchte auch MNR danken für ihre Unterstützung. Wir sind eine kleine Gruppe und weiss nicht, was die Zukunft bringt, aber lasst uns einfach Jesus bezeugen. Das ist unser Auftrag.

    Vielen Dank und Gottes Segen Lothar



\end{document}